% Part: first-order-logic
% Chapter: beyond
% Section: higher-order-logic

\documentclass[../../../include/open-logic-section]{subfiles}

\begin{document}

\olfileid{fol}{byd}{hol}

\olsection{منطق الرتب العليا}

بالانتقال من الرتبة الأولى إلى الثانية صار لنا أن نتناول، في
اللغة الصورية نفسها، مجموعات الأشياء التي في !!{domain}
الرتبة الأولى. ولا موجب للوقوف هنا: ففي الرتبة الثالثة نتناول
مجموعات من مجموعات الأشياء، أو مجموعات تجمع الأشياء ومجموعاتها؛
وفي الرابعة مجموعات من تلك الأنواع. ويتكرر هذا البناء كلما شئنا.

وكثيرًا ما يُصاغ منطق الرتب العليا بالدوال بدل العلاقات، ولا
فرق جوهريًا بين الوجهين بعد إجراء المطابقات الطبيعية. فلنأخذ
«أصنافًا» أساسية $\OLId{latin-looped}{A}$، و$\OLId{latin-looped}{B}$، و$\OLId{latin-looped}{C}$، و~\dots، نسميها من الآن
«أنماطًا»، ولننشئ أنماطًا أخرى بالشرط الآتي:
\begin{quote}
إذا كان $\sigma$ و$\tau$ نمطين منتهيين، كان $\sigma \to \tau$ كذلك.
\end{quote}
الأنماط علامات تركيبية تصنّف الأشياء المقصودة في
!!{domain}نا؛ فالنمط $\sigma \to \tau$ للأشياء التي هي دوال
من أشياء النمط~$\sigma$ إلى أشياء النمط~$\tau$. ومن ذلك أن
نضع نمطًا~$\Omega$ لقيمتي الصدق «صادق» و«كاذب»، ونمطًا~$\Nat$
للأعداد الطبيعية. فعندئذ تكون أشياء النمط $\Nat \to \Omega$
بمثابة العلاقات الأحادية، أي المجموعات الجزئية من~$\Nat$؛
وأشياء النمط $\Nat \to
\Nat$ دوال بين الأعداد الطبيعية؛ وأشياء النمط
$(\Nat \to \Nat) \to \Nat$ «دوال دالية»، أي دوال من نمط أعلى
ترسل الدوال إلى أعداد.

ولنا أن نعتبر منطق الرتب العليا منطقًا متعدد الأصناف، كما
فعلنا في الرتبة الثانية، فنجعل لكل نمط من الأشياء صنفًا.
ولكن الأبين غالبًا أن نبني تركيبه من أصوله. فمن أمثلة ذلك
تعريف الأنماط المنتهية بالاستقراء هكذا:
\begin{enumerate}
\item $\Nat$ نمط منتهٍ.
\item إذا كان $\sigma$ و$\tau$ نمطين منتهيين، كان $\sigma
  \to \tau$ كذلك.
\item إذا كان $\sigma$ و$\tau$ نمطين منتهيين، كان $\sigma \times
  \tau$ كذلك.
\end{enumerate}
يدل $\Nat$، حدسيًا، على نمط الأعداد الطبيعية، ويدل $\sigma \to
\tau$ على نمط الدوال من~$\sigma$ إلى~$\tau$، ويدل $\sigma \times
\tau$ على نمط أزواج الأشياء، أحدهما من~$\sigma$ والآخر من~$\tau$.
ويمكننا عندئذ تعريف مجموعة من الحدود استقرائيًا كما يأتي:
\begin{enumerate}
\item لكل نمط~$\sigma$ مخزون من المتغيرات $\OLId{latin-ordinary}{x}$، و$\OLId{latin-ordinary}{y}$، و$\OLId{latin-ordinary}{z}$، و\dots
  من النمط~$\sigma$.
\item $\Obj 0$ حد من النمط~$\Nat$.
\item $\Obj S$ (الخلف) حد من النمط $\Nat \to \Nat$.
\item إذا كان $\OLId{latin-ordinary}{s}$ حدًا من النمط~$\sigma$، وكان $\OLId{latin-ordinary}{t}$ حدًا من النمط
  $\Nat \to (\sigma \to \sigma)$، كان $\Obj{R}_{st}$ حدًا من النمط
  $\Nat \to \sigma$.
\item إذا كان $\OLId{latin-ordinary}{s}$ حدًا من النمط $\tau \to \sigma$ وكان $\OLId{latin-ordinary}{t}$ حدًا من
  النمط~$\tau$، كان $\OLId{latin-ordinary}{s}(\OLId{latin-ordinary}{t})$ حدًا من النمط~$\sigma$.
\item إذا كان $\OLId{latin-ordinary}{s}$ حدًا من النمط~$\sigma$ وكان $\OLId{latin-ordinary}{x}$ متغيرًا من
  النمط~$\tau$، كان $\lambd[\OLId{latin-ordinary}{x}][\OLId{latin-ordinary}{s}]$ حدًا من النمط $\tau \to \sigma$.
\item إذا كان $\OLId{latin-ordinary}{s}$ حدًا من النمط~$\sigma$ وكان $\OLId{latin-ordinary}{t}$ حدًا من
  النمط~$\tau$، كان $\tuple{\OLId{latin-ordinary}{s}, \OLId{latin-ordinary}{t}}$ حدًا من النمط $\sigma \times
  \tau$.
\item إذا كان $\OLId{latin-ordinary}{s}$ حدًا من النمط~$\sigma \times \tau$، كان $\OLId{latin-ordinary}{p}_\OLMathNumeral{1}(\OLId{latin-ordinary}{s})$
  حدًا من النمط~$\sigma$، و$\OLId{latin-ordinary}{p}_\OLMathNumeral{2}(\OLId{latin-ordinary}{s})$ حدًا من النمط~$\tau$.
\end{enumerate}
وتقريب معنى $\Obj{R}_{st}$ أنه الدالة التي تعرّف عوديًا بالمعادلتين
\begin{align*}
\Obj{R}_{st}(\OLMathNumeral{0}) & = \OLId{latin-ordinary}{s} \\
\Obj{R}_{st}(\OLId{latin-ordinary}{x}+\OLMathNumeral{1}) & = \OLId{latin-ordinary}{t}(\OLId{latin-ordinary}{x}, \OLId{latin-looped}{R}_{st}(\OLId{latin-ordinary}{x})),
\end{align*}
ويدل $\tuple{\OLId{latin-ordinary}{s}, \OLId{latin-ordinary}{t}}$ على الزوج الذي مركبته الأولى~$\OLId{latin-ordinary}{s}$ ومركبته
الثانية~$\OLId{latin-ordinary}{t}$، ويدل $\OLId{latin-ordinary}{p}_\OLMathNumeral{1}(\OLId{latin-ordinary}{s})$ و$\OLId{latin-ordinary}{p}_\OLMathNumeral{2}(\OLId{latin-ordinary}{s})$ على العنصرين الأول والثاني
(«الإسقاطين») من~$\OLId{latin-ordinary}{s}$. وأما $\lambd[\OLId{latin-ordinary}{x}][\OLId{latin-ordinary}{s}]$ فيدل على الدالة~$\OLId{latin-ordinary}{f}$
المعرفة بـ
\[
\OLId{latin-ordinary}{f}(\OLId{latin-ordinary}{x}) = \OLId{latin-ordinary}{s}
\]
لكل~$\OLId{latin-ordinary}{x}$ من النمط~$\tau$؛ ومن ثم يعطينا البند (\OLClassicalDigits{6}) صورة من الاستيعاب،
نعرّف بها الدوال باستعمال الحدود. وتُبنى ال!!{formula}s من
عبارات ال!!{identity}~$\eq[\OLId{latin-ordinary}{s}][\OLId{latin-ordinary}{t}]$ بين حدود من النمط نفسه، والروابط
القضوية المعتادة، والتسوير على الأنماط العليا. فنجعل
المعادلات الأساسية الضابطة لهذه الحدود بديهيات للنسق، ونضمّ إليها
قواعد المنطق المعتادة للمكممات وال!!{identity}.

فإن زدنا على نسق الأنماط المنتهية نمطًا~$\Omega$ لقيم الصدق،
زدنا معه بديهيات تضبط استعماله. بل يمكن، بحسن ترتيب التعريفات،
الاستغناء عن ال!!{formula}s المركبة كلها بحدود من النمط~$\Omega$!{}
ثم يعدّل نسق الإثبات على وفق ذلك، فنحصل في الجوهر على
\emph{نظرية الأنماط البسيطة} التي وضعها ألونزو تشرش في
ثلاثينيات القرن العشرين.

ولدلالة الأنماط العليا وجوه، كما لدلالة الرتبة الثانية.
ففي الدلالة التامة تتراوح متغيرات النمط $\sigma
\to \tau$ على \emph{جميع} الدوال من أشياء النمط~$\sigma$
إلى أشياء النمط~$\tau$. وهذه الدلالة أقوى من أن يكون لها
نسق !!{derivation} فعّال تام. وأما الدلالة الأضعف، فتتألف فيها
!!a{structure} من مجموعات~$\OLId{latin-looped}{T}_\tau$ لكل نمط~$\tau$، ومعها
عمليات التطبيق والإسقاط وغيرها على الوجه المناسب. وبضبط هذه
التفاصيل تثبت مبرهنات تمام لأنساق ال!!{derivation} من الأنواع
الموصوفة آنفًا.

ومن نفع منطق الأنماط العليا أنه يتسع لقدر كبير من الرياضيات
على نحو طبيعي؛ فبدءًا من~$\Nat$ تُعرّف الأعداد الحقيقية والدوال
المتصلة وما إليها. وله مزية خاصة في المنطق الحدسي، إذ تقبل
الأنماط تأويلات بنائية واضحة. ويمكن بناء دلالات للأنماط العليا
على المنطق الحدسي بدل الكلاسيكي، فتتضح بها تلك التأويلات؛ ولها
من وجوه كثيرة فائدة أزيد من نظائرها الكلاسيكية.

\end{document}
